% Polished continuation of Section~\ref{sec:introduction},
% beginning at "Specifically, the framework incorporates:".
% Requires: \usepackage{enumitem}

Specifically, the framework incorporates the following components:

\begin{itemize}[leftmargin=*]

  \item \textbf{Parameterized workload complexity modeling.}
  A classical performance model that projects CPU execution time and memory
  bandwidth contention as a function of qubit count ($2^{n}$ state-space
  scaling), circuit depth, variational parameter count, and the number of
  classical optimization iterations, rather than treating classical
  post-processing as a fixed-cost stage.

  \item \textbf{Unified resource scheduling and brokerage.}
  A broker architecture that manages concurrent, heterogeneous jobs under
  strict capacity constraints. Classical allocation is governed by CPU core
  and memory-bandwidth availability, while quantum allocation is resolved at
  two levels: aggregate qubit capacity and physical connectivity, so that a
  job is admitted only when a connected subgraph of free qubits exists on the
  target device. This exposes fragmentation effects that qubit-count-only
  models cannot capture.

  \item \textbf{Energy and power accounting.}
  A cost model that attributes energy to each execution phase from
  device-specific power profiles, distinguishing the near-constant cryogenic
  baseline of superconducting QPUs from the load-proportional draw of
  classical nodes. Power parameters are grounded in the published consumption
  figures of Ezratty~\cite{ezratty2023}, and are exposed as configuration
  parameters so that alternative hardware assumptions and pricing regimes can
  be evaluated without modifying the simulator.

  \item \textbf{Integrated telemetry and monitoring.}
  An event-driven cloud monitor that records device utilization, instantaneous
  fleet power, and queueing behavior at event boundaries throughout the run,
  producing per-job records and time-series traces. The resulting traces
  reconcile discrete per-job measurements with timeline-integrated,
  system-wide capacity and energy metrics, enabling analysis to be performed
  over the simulation itself rather than over aggregate end-of-run summaries.

\end{itemize}

Through a series of simulation studies, we demonstrate how HybridCloudSim
enables system architects to locate hardware bottlenecks, compare dynamic
scheduling policies, and reason about the energy and cost implications of
resource allocation in emerging hybrid cloud infrastructures. In particular,
we show that energy- and power-aware simulation provides a principled basis on
which researchers and cloud operators can set pricing policy, by varying
device configurations and workload parameters within the framework rather than
by procuring physical hardware. Detailed circuit execution is deliberately
abstracted away, allowing the framework to act as a lightweight digital twin
of a hybrid quantum--classical cloud. As this work targets orchestration-level
modeling, we refer the reader to Nielsen and Chuang~\cite{nielsen2010quantum}
for a comprehensive treatment of the underlying quantum computation and
quantum information theory.
